\begin{abstract}
    Public vulnerability information is distributed across multiple repositories that differ in scope, structure, and update mechanisms, making integrated cybersecurity analyses difficult. An integrated vulnerability intelligence dataset is presented by collecting, normalizing, and correlating data from CVE List V5, the National Vulnerability Database (NVD), the CISA Known Exploited Vulnerabilities Catalog (KEV), and the GitHub Advisory Database. The resulting knowledge graph combines complementary technical and operational attributes describing vulnerabilities, software weaknesses, affected packages, and external references. The complete data construction workflow, quality assurance procedures, and storage model are documented to ensure reproducibility and support future extensions. The resulting resource provides a reusable foundation for vulnerability management, cyber threat intelligence, knowledge graph construction, and machine learning applications.
\end{abstract}

\begin{resumo}
    Informações públicas sobre vulnerabilidades de software encontram-se distribuídas em diferentes repositórios, dificultando análises integradas e a reprodutibilidade de pesquisas em segurança cibernética. É apresentada uma base integrada de inteligência de vulnerabilidades construída a partir da coleta, normalização e integração de dados provenientes da CVE List V5, National Vulnerability Database (NVD), CISA Known Exploited Vulnerabilities (KEV) e GitHub Advisory Database. O resultado consiste em um grafo de conhecimento que reúne atributos técnicos e operacionais sobre vulnerabilidades, fraquezas de software, pacotes afetados e referências externas. O processo de construção, validação e organização dos dados é documentado para favorecer a reprodutibilidade e a evolução contínua da base, oferecendo um recurso reutilizável para gestão de vulnerabilidades, inteligência de ameaças, mineração de dados e aprendizado de máquina.
\end{resumo}